% Work-MAT | 2. rocnik | pripravovana hodina 36
% Návaznost: uživatelské připravy.zip (hodiny 01–20), autorská učebnice na GitHubu.
% Realisticky.cz: didaktická inspirace v posloupnosti a typech úloh; zadání jsou vlastní.
\documentclass[11pt,a4paper]{article}
\input{../../sablony/pripravy/priprava.sty}
\nastavPripravu{2. ročník}{Logaritmické rovnice}{36}{Základní logaritmické rovnice}
\begin{document}
\blok{Motivace a návaznost}
Nejdříve požaduj podmínky logaritmu, potom rovnici převeď na odpovídající exponenciální zápis. Zkouška ověřuje kladný argument.
\poznamka{Orientační nabídka pro 45--55 minut: krátké zopakování a motivace, nový výklad, tři společné vzorové úlohy, samostatná práce. Vyber podle tempa třídy; zbývající zadání slouží jako rezerva.}
\blok{Výklad}
Z rovnice $\log_a F(x)=b$ plyne $F(x)=a^b$, ale pouze pro $F(x)>0$.
\blok{Společné příklady}
\begin{prikladReseniBox}[title={Příklad 1}]
Vyřeš, resp. urči: $\log_3(x-1)=2$.
\tcblower
\[\begin{aligned}x-1>0 \\ x-1=3^2=9 \\ x=10\end{aligned}\]
\end{prikladReseniBox}
\begin{prikladReseniBox}[title={Příklad 2}]
Vyřeš, resp. urči: $\log_2(2x+1)=3$.
\tcblower
\[\begin{aligned}2x+1>0 \\ 2x+1=8\;\Longrightarrow\;x=\frac{7}{2}\end{aligned}\]
\end{prikladReseniBox}
\poznamka{Zadání samostatné práce střídej s výpočtem základních logaritmů; podmínky jsou povinné.}
\blok{Samostatné procvičování}
\begin{zadaniReseniBox}[title={Samostatná práce 1}]
\begin{enumerate}[label=\alph*)]
\item $\log_{2}(x-1)=1$.
\item $\log_{3}(x-2)=2$.
\item $\log_{5}(x-3)=3$.
\item $\log_{2}(x-4)=1$.
\item $\log_{3}(x-1)=2$.
\item $\log_{5}(x-2)=3$.
\end{enumerate}
\tcblower
\begin{enumerate}[label=\alph*)]
\item $x=3$.
\item $x=11$.
\item $x=128$.
\item $x=6$.
\item $x=10$.
\item $x=127$.
\end{enumerate}
\end{zadaniReseniBox}
\blok{Dobrovolné rozšíření}
\begin{zadaniReseniBox}[title={Rozšíření 1}]
\begin{enumerate}[label=\alph*)]
\item $Vytvoř vlastní úlohu tohoto typu a se spolužákem porovnej podmínky a výsledek.$.
\end{enumerate}
\tcblower
\begin{enumerate}[label=\alph*)]
\item Zkontroluj definice a argumenty..
\end{enumerate}
\end{zadaniReseniBox}
\blok{Úkol na doma}
U každého logaritmu vypiš podmínky argumentu a proveď zkoušku. U rovnic vyznač případně nepřípustné kořeny, u nerovnic odpověz intervalem.
\begin{zadaniReseniBox}[title={Domácí úkol}]
\begin{enumerate}[label=\alph*)]
\item $\log_{2}(x-3)=1$.
\item $\log_{3}(x-4)=2$.
\item $\log_{5}(x-1)=3$.
\item $\log_{2}(x-2)=1$.
\item $\log_{3}(x-3)=2$.
\item $\log_{5}(x-4)=3$.
\item $\log_{2}(x-1)=1$.
\item $\log_{3}(x-2)=2$.
\end{enumerate}
\tcblower
\begin{enumerate}[label=\alph*)]
\item $x=5$.
\item $x=13$.
\item $x=126$.
\item $x=4$.
\item $x=12$.
\item $x=129$.
\item $x=3$.
\item $x=11$.
\end{enumerate}
\end{zadaniReseniBox}
\poznamka{Domácí úkol obsahuje 8 položek; předpokládaná domácí práce přibližně 15--30 minut. Po větších blocích jsou vědomě promíchané dříve probrané typy.}
\end{document}
